集合~$A$ 上的一个\emph{关系}~$R$ 是将~$A$ 中元素联系起来的一种方式。
若关系~$R$ 在~$x$ 和~$y$ 之间\emph{成立}，我们就写作 $Rxy$。
形式上，我们可以将~$R$ 视为使得~$Rxy$ 成立的所有有序对
$\tuple{x,y} \in A^2$ 的集合。
小于、大于、等于、整除、长度相等、子集关系以及大小相等
（分别定义在数的集合、字符串的集合或集合族上），这些都是关系的重要例子。
\emph{图}是直观表示关系的一种通用方式。
但一个图也可以看作是一个二元关系（\emph{边}关系）连同其底层的\emph{顶点}集合。

某些关系共有的特性使它们格外有趣或有用。
关系~$R$ 是\emph{自反的}，如果每个元素都与自身有~$R$ 关系；是\emph{对称的}，
如果对任意~$x$ 和~$y$，只要~$Rxy$ 成立就有~$Ryx$ 成立；是\emph{传递的}，
如果~$Rxy$ 和~$Ryz$ 能保证~$Rxz$。
同时具备这三个性质的关系是\emph{等价关系}。
关系~$R$ 是\emph{反对称的}，如果~$Rxy$ 和~$Ryx$ 能保证~$x=y$。
\emph{偏序}就是那些自反、反对称且传递的关系。
\emph{线序}是任意满足下述条件的偏序：
对任意~$x$ 和~$y$，要么~$Rxy$，要么~$x=y$，要么~$Ryx$。
（通常，具有此性质的关系被称为\emph{连通的}）。

由于关系是（有序对的）集合，我们可以像对集合那样对它们进行运算
（例如，可以构造关系的并与交）。我们还可以将它们链接起来
（\emph{相对积}~$R \mid S$）。
如果我们把~$R$ 与其自身作任意多次的相对积，就得到~$R$ 的
\emph{传递闭包}~$R^+$。